\begin{kbvderiv}
\vspace{4pt plus 3pt minus 1pt}
\begin{center}\begin{adjustbox}{max width=\textwidth}
{\renewcommand{\arraystretch}{1.22}\setlength{\tabcolsep}{4pt}
\begin{tabular}{l c l}
ಗಿರಿ + ಇಂದ್ರ & {\color{kbvgrey}$=$} & ಗಿರೀಂದ್ರ .....(ಇ + ಇ = ಈ) \\
ಧ್ವಜ + ಆರೋಹಣ & {\color{kbvgrey}$=$} & ಧ್ವಜಾರೋಹಣ .....(ಅ + ಆ = ಆ) \\
ಲಕ್ಷ್ಮೀ + ಈಶ & {\color{kbvgrey}$=$} & ಲಕ್ಷ್ಮೀಶ .....(ಈ + ಈ = ಈ) \\
\end{tabular}}
\end{adjustbox}\end{center}
\vspace{4pt plus 3pt minus 1pt}
\end{kbvderiv}

\vspace{4pt plus 4pt minus 1pt}\noindent{\kbvsecfont\fontsize{11pt}{14.5pt}\selectfont\color{kbvaccentdark} (೨)\hspace{0.45em}ಗುಣ ಸಂಧಿ :}\hspace{0.5em}“ಗಣೇಶ” ಈ ಶಬ್ದವನ್ನು ಬಿಡಿಸಿ ನೋಡೋಣ.\par

\begin{kbvderiv}
\vspace{4pt plus 3pt minus 1pt}
\begin{center}\begin{adjustbox}{max width=\textwidth}
{\renewcommand{\arraystretch}{1.22}\setlength{\tabcolsep}{4pt}
\begin{tabular}{l c l}
ಗಣೇಶ & {\color{kbvgrey}$=$} & ಗಣ + ಈಶ \\
 & {\color{kbvgrey}$=$} & ಗ + ಣ್‌ + (ಅ + ಈ) + ಶ \\
\end{tabular}}
\end{adjustbox}\end{center}
\vspace{4pt plus 3pt minus 1pt}
\end{kbvderiv}

\noindent\hspace*{2.2em}ಇಲ್ಲಿ ಅ ಕಾರಕ್ಕೆ ಈ ಕಾರ ಪರವಾಗಿದೆ.\par

\indent  = ಗ + ಣ್‌ + (ಏ) + ಶ ಇಲ್ಲಿ ಅ, ಈ ಗಳ ಸ್ಥಾನದಲ್ಲಿ ಏ ಕಾರ ಆದೇಶವಾಗಿ ಬಂದಿದೆ.\par

\indent  = ಗಣೇಶ .....(\ensuremath{\because} ಣ್‌ + ಈ = ಣೇ) ತೀರ್ಮಾನ : ಅ ಕಾರಕ್ಕೆ ಈ ಕಾರ ಪರವಾದಾಗ ಏ ಕಾರ ಆದೇಶವಾಗಿ ಬರುವುದು.\par

\indent ಇದೇ ರೀತಿಯಾಗಿ,\par

\begin{kbvderiv}
\vspace{4pt plus 3pt minus 1pt}
\begin{center}\begin{adjustbox}{max width=\textwidth}
{\renewcommand{\arraystretch}{1.22}\setlength{\tabcolsep}{4pt}
\begin{tabular}{l c l}
ದೇವ + ಇಂದ್ರ & {\color{kbvgrey}$=$} & ದೇವೇಂದ್ರ .....(ಅ + ಇ = ಏ) \\
ಮಹಾ + ಈಶ್ವರೀ & {\color{kbvgrey}$=$} & ಮಹೇಶ್ವರೀ .....(ಆ + ಈ = ಏ) \\
ಜನ + ಉಪಕಾರಿ & {\color{kbvgrey}$=$} & ಜನೋಪಕಾರಿ .....(ಅ + ಉ = ಓ) \\
ಮಹಾ + ಉನ್ನತಿ & {\color{kbvgrey}$=$} & ಮಹೋನ್ನತಿ .....(ಅ + ಉ = ಓ) \\
ರಾಜ + ಋಷಿ & {\color{kbvgrey}$=$} & ರಾಜರ್ಷಿ .....(ಅ + ಋ = ಅರ್‌) \\
ದೇವ + ಋಷಿ & {\color{kbvgrey}$=$} & ದೇವರ್ಷಿ .....(ಅ + ಋ = ಅರ್‌) \\
\end{tabular}}
\end{adjustbox}\end{center}
\vspace{4pt plus 3pt minus 1pt}
\end{kbvderiv}

\indent ಈ ಉದಾಹರಣೆಗಳಲ್ಲಿ ಪೂರ್ವದ ಸ್ವರವು “ಅ” ಅಥವಾ “ಆ” ಆಗಿರುವುದು ; ಅವುಗಳಿಗೆ ಪರವಾಗಿರುವ ಸ್ವರಗಳು ಇ, ಈ, ಉ, ಋ ಆಗಿರುವುವು. ಅ, ಆಗಳಿಗೆ ಇ, ಈ ಗಳು ಪರವಾದಾಗ ಏಕಾರವೂ, ಅಕಾರಕ್ಕೆ ಉಕಾರ ಪರವಾದಾಗ ಓಕಾರವೂ, ಅಕಾರಕ್ಕೆ ಋಕಾರ ಪರವಾದಾಗ “ಅರ್‌” ಕಾರವೂ, ಆದೇಶಗಳಾಗಿ ಬಂದಿರುವುವು. ಈ ತೆರನಾದ ಸಂಧಿಯು ಗುಣ ಸಂಧಿಯೆನಿಸುವುದು.\par
